\subsection{格}

定义8. 一个有序集E被称为格，如果E的每个由两个元素组成的子集在E中都有最小上界和最大下界。

格的每个乘积都是格；这由有序集乘积中最小上界存在的条件（第9号，命题8）得出。集合A的子集所构成的集合，按包含关系排序，是一个格，因为A的两个子集的并集和交集仍然是A的子集。

\minorheading{示例}

* (1) 整数集 \(\geqslant 1\) ，按 \(m\) 与 \(n\) 之间的关系“ \(m\) 整除 \(n\) ”排序，是一个格； \(\{m, n\}\) 的最小上界是 \(m\) 与 \(n\) 的最小公倍数，而最大下界是它们的最大公因数。

\begin{enumerate}
\def\labelenumi{(\arabic{enumi})}
\setcounter{enumi}{1}
\item
  群G的子群集合，按包含关系排序，构成一个格。
\item
  集合A上的所有拓扑所成的集合，按 \(\mathfrak{G}\) 与 \(\mathfrak{G}'\) 之间的关系“ \(\mathfrak{G}\) 比 \(\mathfrak{G}'\) 更粗”排序，是一个格。（一般拓扑学，第一章，§2）。
\item
  集合 \(\mathcal{F}(\mathbf{I},\mathbf{R})\) ，即定义在区间 \(\mathbf{I}\) （ \(\mathbf{R}\) 中的区间）上的所有实值函数所成的集合，关于序关系 \(f\leqslant g\) （第4号）是一个格，并且作为这样的格同构于乘积 \(\mathbf{R}^{\mathrm{I}}\) . *
\end{enumerate}

注记。一个格显然既是左有向的又是右有向的。但一个既是左有向又是右有向的有序集不一定是格。* 后者的一个例子是映射 \(x \to p(x)\) （从 \(\mathbf{R}\) 到其自身的映射）所成的集合，其中 \(p\) 是 \(\mathbf{R}[\mathbf{X}]\) 中的多项式，该集合按关系 \(p \leqslant q\) 排序（第4号）。*

